% =====================================================================
% Chapter 6: The plane-wave condition as a general tool
% Theory session, dispersion branch (Nov 2026)
% Requires: amsmath, amssymb
% Include in the main file with  % =====================================================================
% Chapter 6: The plane-wave condition as a general tool
% Theory session, dispersion branch (Nov 2026)
% Requires: amsmath, amssymb
% Include in the main file with  % =====================================================================
% Chapter 6: The plane-wave condition as a general tool
% Theory session, dispersion branch (Nov 2026)
% Requires: amsmath, amssymb
% Include in the main file with  % =====================================================================
% Chapter 6: The plane-wave condition as a general tool
% Theory session, dispersion branch (Nov 2026)
% Requires: amsmath, amssymb
% Include in the main file with  \include{chapter_06_pwe_general}
% =====================================================================

\chapter{The plane-wave condition as a general tool}
\label{ch:pwe_general}

This chapter takes the plane-wave expansion (PWE) condition of
Chapter~\ref{ch:gap_position} and asks what can be read off it without solving
it: how many bands exist, where they can lie, what happens to them when the
resonator is detuned, and what changes with more than one resonator per cell.
Numbers come from the Reissner--Mindlin (RM) plane-wave model with $\mu = 0.25$
and $f_r = 1200\,\mathrm{Hz}$ unless stated otherwise.

\begin{center}
\begin{tabular}{|l|p{0.62\textwidth}|}
\hline
Symbol & Meaning \\
\hline
$\lambda = \omega^2$ & eigenvalue of one dot in the diagram \\
$\lambda_n^0$ & bare plate eigenvalue of wave $n$ at this $\mathbf{k}$ \\
$\lambda_r = k_s/m$ & the resonator's own eigenvalue ($k_s$ spring, $m$ mass) \\
$\mu = m/(\rho h A)$ & resonator mass over plate mass per cell ($\rho h$ mass per area, $A$ cell area) \\
$\beta = \lambda/\lambda_r$ & squared frequency relative to the resonator \\
$w_0$ & plate deflection at the resonator \\
$Z$, $C$ & force per deflection of the resonator; deflection per force of the plate \\
$P(\lambda)$ & plate curve, Eq.~\eqref{eq:plate_curve} \\
$Y(\lambda)$ & resonator line, Eq.~\eqref{eq:resonator_line} \\
$S$ & at $\Gamma$: the part of $P$ from all waves except $\mathbf{G}=0$ \\
$g(\kappa, \omega)$ & deflection per unit point force of one plane wave of wavenumber $\kappa$ \\
$D$, $D_s$ & bending stiffness and shear stiffness of the plate \\
$I_r = \rho h^3/12$ & rotary inertia per area \\
$f_X$ & bare plate frequency at $X$, $1194.4\,\mathrm{Hz}$ \\
\hline
\end{tabular}
\end{center}

% ---------------------------------------------------------------------
\section{Recap of the condition}

Three facts combine:
\begin{itemize}
  \item A resonator exerts a force on the plate like a mass that depends on
        frequency: $F = m_\mathrm{eff}\,\lambda\, w_0$, with
        $m_\mathrm{eff} = m/(1-\beta)$.
  \item A point force excites every mass-normalized bare wave $\varphi_n$ with
        amplitude
        \[ q_n = \frac{F\,\varphi_n(\mathbf{x}_0)^*}{\lambda_n^0 - \lambda} . \]
  \item Adding up all waves at $\mathbf{x}_0$ must give back $w_0$.
\end{itemize}
Inserting one into the other, with
$|\varphi_n(\mathbf{x}_0)|^2 = 1/(\rho h A)$ and
$\mu_\mathrm{eff} = \mu/(1-\beta)$, gave
\begin{equation}
  1 = \mu_\mathrm{eff}\,\sum_n \frac{\lambda}{\lambda_n^0 - \lambda} .
  \label{eq:pwe_recap}
\end{equation}

% ---------------------------------------------------------------------
\section{Separating resonator and plate}

Name the sum:
\begin{equation}
  P(\lambda) = \sum_n \frac{\lambda}{\lambda_n^0 - \lambda} .
  \label{eq:plate_curve}
\end{equation}
Then Eq.~\eqref{eq:pwe_recap} reads $1 = \mu P/(1-\beta)$. Multiply both sides
by $(1-\beta)/\mu$:
\begin{equation}
  \frac{1-\beta}{\mu} = P(\lambda) .
\end{equation}
Name the left side:
\begin{equation}
  Y(\lambda) = \frac{1 - \lambda/\lambda_r}{\mu} .
  \label{eq:resonator_line}
\end{equation}
The condition is then simply
\begin{equation}
  \boxed{\;Y(\lambda) = P(\lambda)\;}
  \label{eq:Y_equals_P}
\end{equation}
The two sides have separate origins:
\begin{itemize}
  \item $Y$ contains only the resonator ($\lambda_r$, $\mu$). It is a falling
        straight line in $\lambda$, with $Y = 1/\mu$ at $\lambda = 0$ and
        $Y = 0$ at $\lambda = \lambda_r$.
  \item $P$ contains only the plate and $\mathbf{k}$. It does not know about
        the resonator.
\end{itemize}

\paragraph{$P$ rises everywhere.} With the quotient rule, for one term,
\begin{equation}
  \frac{\mathrm{d}}{\mathrm{d}\lambda}\,\frac{\lambda}{\lambda_n^0 - \lambda}
  = \frac{(\lambda_n^0 - \lambda) + \lambda}{(\lambda_n^0 - \lambda)^2}
  = \frac{\lambda_n^0}{(\lambda_n^0 - \lambda)^2} \ge 0 .
\end{equation}

\paragraph{$P$ has a pole at every distinct bare value.}
Just below $\lambda_n^0$ the denominator is small and positive, so
$P \to +\infty$. Just above it, $P \to -\infty$. Between two neighbouring bare
values, $P$ therefore climbs from $-\infty$ to $+\infty$.

The dots of the diagram at this $\mathbf{k}$ are the intersections of a falling
line with these rising branches.

% ---------------------------------------------------------------------
\section{What follows from that picture}

Consider $Y - P$ between two neighbouring bare values. It is a falling function
minus a rising one, so it falls strictly, from $+\infty$ to $-\infty$, and
crosses zero \emph{exactly once}. The same holds below the lowest bare value,
where $Y - P$ starts positive.

\paragraph{Consequence 1: bands are trapped.}
There is exactly one root between neighbouring bare frequencies. A single point
resonator cannot push a band past a bare frequency. Example at $X$:
\begin{center}
\begin{tabular}{|p{0.5\textwidth}|p{0.33\textwidth}|}
\hline
Interval & Root \\
\hline
below $1194.4\,\mathrm{Hz}$ & $835.6\,\mathrm{Hz}$ \\
exactly at $1194.4\,\mathrm{Hz}$ & sin pattern (node at the resonator) \\
between $1194.4\,\mathrm{Hz}$ and the next bare value & $1652.4\,\mathrm{Hz}$ \\
\hline
\end{tabular}
\end{center}

\paragraph{Consequence 2: the stuck modes.}
A $d$-fold bare value is one pole, so it supplies only one interval boundary. The
remaining $d-1$ combinations have $w_0 = 0$ and sit exactly on the pole. Counting
all roots gives all bare modes plus one extra root per $\mathbf{k}$: the extra
band of the one extra resonator degree of freedom.

\paragraph{Consequence 3: which side of $f_r$ a root is on.}
At a root, $Y = P$, and $Y > 0$ exactly when $\lambda < \lambda_r$. Therefore:
\begin{itemize}
  \item $P > 0$ at the root: the root is \emph{below} $f_r$; the plate at the
        attachment acts spring-like overall;
  \item $P < 0$ at the root: the root is \emph{above} $f_r$; the plate at the
        attachment acts mass-like overall.
\end{itemize}
Below the lowest bare value every term is positive, so $P > 0$ and band 1 lies
below $f_r$. At $\Gamma$, $P = -1 + S$, where $-1$ is the $\mathbf{G}=0$ wave
(rigid cell) and $S$ is the sum over all others. The band at $\Gamma$ sits above
$f_r$ only while $S < 1$.

% ---------------------------------------------------------------------
\section{The sum for Reissner--Mindlin}

In the code (\texttt{plane\_wave\_point\_compliance}) each wave enters through
$g(\kappa, \omega)$. For Kirchhoff,
\begin{equation}
  g = \frac{1}{D\kappa^4 - \rho h\,\omega^2} .
\end{equation}
For RM a wave has two coupled unknowns, the deflection and the rotation along
the wave, with dynamic stiffness matrix
\begin{equation}
  \begin{pmatrix}
    D_s\kappa^2 - \rho h\,\omega^2 & -D_s\kappa \\
    -D_s\kappa & D\kappa^2 + D_s - I_r\omega^2
  \end{pmatrix}
  =
  \begin{pmatrix} a & b \\ b & c \end{pmatrix} ,
\end{equation}
and $g$ is the deflection entry of its inverse, $g = c/(ac - b^2)$. At large
$\kappa$ the inertia terms are negligible:
\begin{align}
  ac - b^2 &= D_s\kappa^2\,(D\kappa^2 + D_s) - D_s^2\kappa^2 = D\,D_s\,\kappa^4 , \\
  g &\approx \frac{D\kappa^2 + D_s}{D\,D_s\,\kappa^4}
     = \frac{1}{D\kappa^4} + \frac{1}{D_s\kappa^2} .
\end{align}
Bending and shear compliance simply add. Kirchhoff is the limit
$D_s \to \infty$.

\paragraph{Why this matters for the sum.}
In 2D the number of waves with wavenumber between $\kappa$ and $\kappa +
\mathrm{d}\kappa$ is proportional to the area of the ring, $\kappa\,\mathrm{d}\kappa$.
The tail of the sum therefore behaves like
\begin{itemize}
  \item Kirchhoff: $\int \kappa\,\mathrm{d}\kappa/\kappa^4$, which \emph{converges};
  \item RM shear term: $\int \kappa\,\mathrm{d}\kappa/\kappa^2 = \int \mathrm{d}\kappa/\kappa = \ln\kappa$,
        which \emph{diverges slowly}.
\end{itemize}
Physically, a true point force on a shear-deformable plate gives infinite local
deflection. Real attachments have a finite size; the cutoff in the sum, or the
mesh in FE, plays that role.
\begin{center}
\begin{tabular}{|l|l|l|l|}
\hline
Cutoff & $\Gamma$ band 2, RM & $M$ band 1, RM & $\Gamma$ band 2, Kirchhoff \\
\hline
6 & 1276.1 & 1033.1 & 1280.2 \\
12 & 1274.7 & 1032.1 & 1279.8 \\
24 & 1273.5 & 1031.4 & 1279.7 \\
48 & 1272.4 & 1030.7 & 1279.6 \\
96 & 1271.3 & 1030.0 & -- \\
\hline
\end{tabular}
\end{center}
All values are in Hz; ``cutoff'' is the largest harmonic index kept in the sum.
Each doubling lowers the RM values by about the same amount
($\approx 1.1\,\mathrm{Hz}$ at $\Gamma$, $\approx 0.7\,\mathrm{Hz}$ at $M$), the
signature of a logarithm. The Kirchhoff values settle. FE shows the same effect
under mesh refinement ($1273.5$ / $1271.7$ / $1270.5\,\mathrm{Hz}$ for $10$/$20$/$40$
elements). It is about $0.1\,\%$ per doubling, so small in practice. The upper
gap edge at $X$ is not affected, because the sin pattern does not feel the point
force.

% ---------------------------------------------------------------------
\section{More resonators per cell}

Take resonators at positions $\mathbf{x}_j$, with forces $F_j$ and plate
deflections $w_j$. The steps are the same as for one resonator:
\begin{enumerate}
  \item Each wave feels all the forces:
        $(\lambda_n^0 - \lambda)\, q_n = \sum_l F_l\,\varphi_n(\mathbf{x}_l)^*$.
  \item Solve for $q_n$ and insert into
        $w_j = \sum_n \varphi_n(\mathbf{x}_j)\, q_n$:
        \begin{equation}
          w_j = \sum_l C_{jl}\, F_l ,
          \qquad
          C_{jl} = \sum_n \frac{\varphi_n(\mathbf{x}_j)\,\varphi_n(\mathbf{x}_l)^*}{\lambda_n^0 - \lambda} .
        \end{equation}
  \item Each resonator gives $F_l = Z_l\, w_l$. In matrix form $w = C Z w$, with
        $Z$ diagonal, so a nonzero solution needs
        \begin{equation}
          \boxed{\;\det(I - C Z) = 0\;}
        \end{equation}
\end{enumerate}
With one resonator this is $1 = ZC$. $C_{jl}$ is the deflection at resonator $j$
per unit force at resonator $l$. Each added resonator adds one band and can
couple to patterns that the others cannot.

\paragraph{Where the plane-wave model stops being valid.}
It assumes a uniform plate and point attachments. Cutouts, stiffeners or real
attachment patches need FE. Inside its validity it is nearly free to evaluate, so
it is both the check on FE and a fast design tool.

% ---------------------------------------------------------------------
\section{How detuning moves the edges}

\paragraph{Raising $f_r$.}
$Y = (1 - \lambda/\lambda_r)/\mu$ keeps its value $1/\mu$ at $\lambda = 0$, while
its zero moves to the right, so $Y$ increases at every $\lambda > 0$. At the old
root $Y - P$ is now positive, and since $Y - P$ falls with $\lambda$, the root
moves right: \emph{every coupled root rises with $f_r$}.

\paragraph{Raising $\mu$.}
$Y$ is zero at $\lambda = \lambda_r$ for any $\mu$, so raising $\mu$ turns the
line about that point. Below $\lambda_r$, $Y$ is positive and shrinks, so roots
below $f_r$ go \emph{down}. Above $\lambda_r$, $Y$ is negative and gets closer to
zero, so roots above $f_r$ go \emph{up}. A heavier resonator pushes roots
\emph{away from} $f_r$, which is level repulsion. Check at $f_r = 1200\,\mathrm{Hz}$,
$\mu = 0.25 \to 0.5$: band 3 at $X$ goes from $1653$ to $1842\,\mathrm{Hz}$, band 1
at $M$ from $1032$ to $927\,\mathrm{Hz}$.

\begin{center}
\begin{tabular}{|l|l|l|l|}
\hline
$f_r$ [Hz] & $M$ band 1 [Hz] & $\Gamma$ band 2 [Hz] & $\Gamma$ band 2 $/\, f_r$ \\
\hline
300 & 297 & 334 & 1.11 \\
800 & 748 & 874 & 1.09 \\
1200 & 1031 & 1273 & 1.06 \\
1500 & 1188 & 1546 & 1.03 \\
2000 & 1360 & 1940 & 0.97 \\
5000 & 1600 & 2957 & 0.59 \\
200000 & 1648 & 3279 & 0.02 \\
\hline
\end{tabular}
\end{center}

\paragraph{Limits.}
\begin{itemize}
  \item $f_r \to 0$: $Y$ falls very steeply, so the root sits at
        $\lambda \approx \lambda_r$, and band 1 at $M$ follows $f_r$.
  \item $f_r \to \infty$: $Y$ becomes the horizontal line $1/\mu$, so
        Eq.~\eqref{eq:Y_equals_P} is the rigid-mass condition. The spring is
        infinitely stiff, and band 1 at $M$ saturates at $1648\,\mathrm{Hz}$. The
        four-wave rigid-mass formula $2378/\sqrt{1+4\mu}$ gives
        $1682\,\mathrm{Hz}$; the other waves account for the rest.
\end{itemize}

\paragraph{Three tuning values mark the changes.}
\begin{center}
\begin{tabular}{|l|p{0.62\textwidth}|}
\hline
$f_r$ & What happens \\
\hline
$\approx 1118\,\mathrm{Hz}$ & $\Gamma$ band 2 reaches $f_X$, and the upper edge becomes pinned \\
$\approx 1513\,\mathrm{Hz}$ & $M$ band 1 reaches $f_X$, and the gap closes \\
$\approx 1765\,\mathrm{Hz}$ & $\Gamma$ band 2 equals $f_r$: $S$ reaches $1$, and the plate at $\Gamma$ turns from mass-like to spring-like \\
\hline
\end{tabular}
\end{center}
\paragraph{Why the gap is widest near $1118\,\mathrm{Hz}$.}
Below it, both edges move up with $f_r$, the upper one faster, so the width
grows. Above it, the upper edge is frozen while the lower one keeps rising, so
the width shrinks.

% ---------------------------------------------------------------------
\section{Two resonators per cell: numbers}
\label{sec:two_resonators}

Identical resonators, $f_r = 1200\,\mathrm{Hz}$, plane-wave model on
$\Gamma$--$X$--$M$--$Y$--$\Gamma$ with $15$ $\mathbf{k}$-points per segment. The
model reproduces the one-resonator values to about $1\,\mathrm{Hz}$, but it is not
FE. The path needs $Y$ because a resonator at $(a/2, 0)$ breaks the square
symmetry. The two-resonator gaps lie between bands 2 and 3, because each resonator
adds a band.
\begin{center}
\begin{tabular}{|p{0.23\textwidth}|p{0.15\textwidth}|p{0.18\textwidth}|l|l|}
\hline
Resonators per cell & At $X$ [Hz] & At $Y$ [Hz] & Gap [Hz] & Width \\
\hline
one at $(0,0)$, $\mu = 0.25$ & 835, \textbf{1194}, 1653 & 835, \textbf{1194}, 1653 & 1031--1194 & 163 \\
one at $(0,0)$, $\mu = 0.5$ & 724, \textbf{1194}, 1841 & 724, \textbf{1194}, 1841 & 926--1194 & 269 \\
$(0,0) + (a/2, 0)$, $\mu = 0.25$ each & 835 $\times 2$, 1653 $\times 2$ & 735, 1137, \textbf{1194}, 1913 & 1137--1194 & 58 \\
$(0,0) + (a/2, a/2)$, $\mu = 0.25$ each & 835 $\times 2$, 1653 $\times 2$ & 835 $\times 2$, 1653 $\times 2$ & 1114--1429 & 315 \\
\hline
\end{tabular}
\end{center}

\paragraph{The diagonal arrangement is a smaller lattice.}
Two identical resonators at $(0,0)$ and $(a/2, a/2)$ form a square lattice rotated
by $45^\circ$ with spacing $a/\sqrt{2}$. The first Bragg frequency moves up, from
$X$ at $1194\,\mathrm{Hz}$ to the old $M$ at $2378\,\mathrm{Hz}$. A second resonator
and a smaller cell are therefore the same physics here. The gap then contains
$f_r$, so it is a local-resonance gap.

\paragraph{Extra question: why is the $(0,0)+(a/2,0)$ gap so narrow?}
Sorted values at $Y$:
\begin{center}
\begin{tabular}{|l|l|p{0.58\textwidth}|}
\hline
Band at $Y$ & $f$ [Hz] & Mode \\
\hline
1 & 735 & $\cos(\pi y/a)$, both resonators moving together, lower branch \\
2 & \textbf{1137} & the two resonators moving in antiphase, new mode \\
3 & \textbf{1194.4} & $\sin(\pi y/a)$, node at both resonators, bare \\
4 & 1913 & $\cos(\pi y/a)$, both resonators moving together, upper branch \\
\hline
\end{tabular}
\end{center}
$1137\,\mathrm{Hz}$ is the \emph{maximum} of band 2, so it is the new lower edge;
the upper edge is the minimum of band 3, the pinned mode. The gap shrinks because
the second resonator adds a band, and at $Y$ that band lands inside the old gap.

Why it sits at $1137\,\mathrm{Hz}$: both resonators have $y = 0$, so the waves at
$Y$, which vary only in $y$, have the same value at both. If the resonators move
in antiphase, $F_1 = -F_2$, their forces on those waves cancel exactly, and the $Y$
waves stay at rest. The opposite forces do excite waves that alternate in $x$
between the two resonators: $(\pm 2\pi/a, \pm\pi/a)$, at about
$5868\,\mathrm{Hz}$, and higher waves with odd multiples of $2\pi/a$ in $x$. For
such a wave the two resonators contribute with opposite sign and the antiphase
forces add up, giving a factor $2$:
\begin{equation}
  1 = 2\,\mu_\mathrm{eff}\, P_\mathrm{odd},
  \qquad
  P_\mathrm{odd} = \sum_{\text{odd waves}} \frac{\lambda}{\lambda_n^0 - \lambda} .
\end{equation}
Insert $\mu_\mathrm{eff} = \mu/(1-\beta)$ and multiply by $(1-\beta)$:
\begin{equation}
  1 - \beta = 2\mu\, P_\mathrm{odd} .
\end{equation}
All of these waves lie far above $1137\,\mathrm{Hz}$, so $P_\mathrm{odd} > 0$ and
the plate acts like a soft spring, which pulls the mode below $f_r$. The four waves
at $5868\,\mathrm{Hz}$ give $P_\mathrm{odd} \approx 0.16$, and the higher waves
bring it to about $0.2$, so
\begin{equation}
  \beta \approx 1 - 2\cdot 0.25\cdot 0.205 = 0.897 ,
  \qquad
  f \approx 1200\cdot\sqrt{0.897} \approx 1137\,\mathrm{Hz}.
\end{equation}

% ---------------------------------------------------------------------
\section{Check questions with answers}

\paragraph{Question 1.} \emph{As $f_r$ goes up, band 1 at $M$ rises but saturates
near $1650\,\mathrm{Hz}$, well below its bare value of $2378\,\mathrm{Hz}$. Why does
it stop there, and what has the resonator physically become?}

\emph{Answer.} With $f_r \to \infty$ the spring cannot deform, so the resonator is
a rigid mass glued to the plate. Band 1 at $M$ settles at the mass-loaded value:
$1682\,\mathrm{Hz}$ from the four-wave formula, $1648\,\mathrm{Hz}$ with all waves.
A rigid mass can only lower a frequency, so stiffening the spring further changes
nothing.

\paragraph{Question 2.} \emph{A heavier resonator pushes roots away from $f_r$.
Where in the $Y$--$P$ picture do you see this, and what does it mean for a gap that
contains $f_r$ compared with one lying entirely below $f_r$?}

\emph{Answer.} Raising $\mu$ turns $Y$ about its zero at $\lambda = \lambda_r$.
Roots left of that point move left, roots right of it move right.
\begin{center}
\begin{tabular}{|p{0.25\textwidth}|l|l|p{0.17\textwidth}|}
\hline
Tuning & Gap at $\mu = 0.25$ & Gap at $\mu = 0.5$ & What moved \\
\hline
$f_r = 800\,\mathrm{Hz}$ (contains $f_r$) & 748--874\,Hz & 707--936\,Hz & both edges, in opposite directions \\
$f_r = 1200\,\mathrm{Hz}$ (below $f_r$) & 1032--1194\,Hz & 927--1194\,Hz & lower edge only \\
\hline
\end{tabular}
\end{center}
Widening does not make the $1200\,\mathrm{Hz}$ gap a local-resonance gap. A gap that
contains $f_r$ has one edge on each side of the pivot and is a true local-resonance
gap. The $1200\,\mathrm{Hz}$ gap lies entirely below it, and its upper edge is the
sin pattern at $X$, which is not a root of $Y = P$ at all because it has a node at
the resonator. It stays a Bragg gap with a resonance-boosted mass.

\paragraph{Question 3.} \emph{Add a second resonator in each cell at $(a/2, 0)$, same
$f_r$. What happens to the sin pattern at $X$, and would you expect the upper edge to
stay pinned at $1194\,\mathrm{Hz}$?}

\emph{Answer.} At $X$, $\cos(\pi x/a)$ has its antinode at the first resonator and a
node at the second, and $\sin(\pi x/a)$ the reverse. That makes two identical
plate-plus-resonator pairs. The sin pattern does not simply rise; it \emph{splits}
into a lower and an upper branch, like cos, and the $1194.4\,\mathrm{Hz}$ mode is
gone from $X$.

But $X$ is not the only point that matters. At $Y$ the pattern $\sin(\pi y/a)$ is
zero at $y = 0$, where both resonators sit, so it stays at $1194.4\,\mathrm{Hz}$ and
the upper edge stays pinned. Moving the second resonator to the cell centre removes
the pinning at both $X$ and $Y$ (Section~\ref{sec:two_resonators}).

% =====================================================================

\chapter{The plane-wave condition as a general tool}
\label{ch:pwe_general}

This chapter takes the plane-wave expansion (PWE) condition of
Chapter~\ref{ch:gap_position} and asks what can be read off it without solving
it: how many bands exist, where they can lie, what happens to them when the
resonator is detuned, and what changes with more than one resonator per cell.
Numbers come from the Reissner--Mindlin (RM) plane-wave model with $\mu = 0.25$
and $f_r = 1200\,\mathrm{Hz}$ unless stated otherwise.

\begin{center}
\begin{tabular}{|l|p{0.62\textwidth}|}
\hline
Symbol & Meaning \\
\hline
$\lambda = \omega^2$ & eigenvalue of one dot in the diagram \\
$\lambda_n^0$ & bare plate eigenvalue of wave $n$ at this $\mathbf{k}$ \\
$\lambda_r = k_s/m$ & the resonator's own eigenvalue ($k_s$ spring, $m$ mass) \\
$\mu = m/(\rho h A)$ & resonator mass over plate mass per cell ($\rho h$ mass per area, $A$ cell area) \\
$\beta = \lambda/\lambda_r$ & squared frequency relative to the resonator \\
$w_0$ & plate deflection at the resonator \\
$Z$, $C$ & force per deflection of the resonator; deflection per force of the plate \\
$P(\lambda)$ & plate curve, Eq.~\eqref{eq:plate_curve} \\
$Y(\lambda)$ & resonator line, Eq.~\eqref{eq:resonator_line} \\
$S$ & at $\Gamma$: the part of $P$ from all waves except $\mathbf{G}=0$ \\
$g(\kappa, \omega)$ & deflection per unit point force of one plane wave of wavenumber $\kappa$ \\
$D$, $D_s$ & bending stiffness and shear stiffness of the plate \\
$I_r = \rho h^3/12$ & rotary inertia per area \\
$f_X$ & bare plate frequency at $X$, $1194.4\,\mathrm{Hz}$ \\
\hline
\end{tabular}
\end{center}

% ---------------------------------------------------------------------
\section{Recap of the condition}

Three facts combine:
\begin{itemize}
  \item A resonator exerts a force on the plate like a mass that depends on
        frequency: $F = m_\mathrm{eff}\,\lambda\, w_0$, with
        $m_\mathrm{eff} = m/(1-\beta)$.
  \item A point force excites every mass-normalized bare wave $\varphi_n$ with
        amplitude
        \[ q_n = \frac{F\,\varphi_n(\mathbf{x}_0)^*}{\lambda_n^0 - \lambda} . \]
  \item Adding up all waves at $\mathbf{x}_0$ must give back $w_0$.
\end{itemize}
Inserting one into the other, with
$|\varphi_n(\mathbf{x}_0)|^2 = 1/(\rho h A)$ and
$\mu_\mathrm{eff} = \mu/(1-\beta)$, gave
\begin{equation}
  1 = \mu_\mathrm{eff}\,\sum_n \frac{\lambda}{\lambda_n^0 - \lambda} .
  \label{eq:pwe_recap}
\end{equation}

% ---------------------------------------------------------------------
\section{Separating resonator and plate}

Name the sum:
\begin{equation}
  P(\lambda) = \sum_n \frac{\lambda}{\lambda_n^0 - \lambda} .
  \label{eq:plate_curve}
\end{equation}
Then Eq.~\eqref{eq:pwe_recap} reads $1 = \mu P/(1-\beta)$. Multiply both sides
by $(1-\beta)/\mu$:
\begin{equation}
  \frac{1-\beta}{\mu} = P(\lambda) .
\end{equation}
Name the left side:
\begin{equation}
  Y(\lambda) = \frac{1 - \lambda/\lambda_r}{\mu} .
  \label{eq:resonator_line}
\end{equation}
The condition is then simply
\begin{equation}
  \boxed{\;Y(\lambda) = P(\lambda)\;}
  \label{eq:Y_equals_P}
\end{equation}
The two sides have separate origins:
\begin{itemize}
  \item $Y$ contains only the resonator ($\lambda_r$, $\mu$). It is a falling
        straight line in $\lambda$, with $Y = 1/\mu$ at $\lambda = 0$ and
        $Y = 0$ at $\lambda = \lambda_r$.
  \item $P$ contains only the plate and $\mathbf{k}$. It does not know about
        the resonator.
\end{itemize}

\paragraph{$P$ rises everywhere.} With the quotient rule, for one term,
\begin{equation}
  \frac{\mathrm{d}}{\mathrm{d}\lambda}\,\frac{\lambda}{\lambda_n^0 - \lambda}
  = \frac{(\lambda_n^0 - \lambda) + \lambda}{(\lambda_n^0 - \lambda)^2}
  = \frac{\lambda_n^0}{(\lambda_n^0 - \lambda)^2} \ge 0 .
\end{equation}

\paragraph{$P$ has a pole at every distinct bare value.}
Just below $\lambda_n^0$ the denominator is small and positive, so
$P \to +\infty$. Just above it, $P \to -\infty$. Between two neighbouring bare
values, $P$ therefore climbs from $-\infty$ to $+\infty$.

The dots of the diagram at this $\mathbf{k}$ are the intersections of a falling
line with these rising branches.

% ---------------------------------------------------------------------
\section{What follows from that picture}

Consider $Y - P$ between two neighbouring bare values. It is a falling function
minus a rising one, so it falls strictly, from $+\infty$ to $-\infty$, and
crosses zero \emph{exactly once}. The same holds below the lowest bare value,
where $Y - P$ starts positive.

\paragraph{Consequence 1: bands are trapped.}
There is exactly one root between neighbouring bare frequencies. A single point
resonator cannot push a band past a bare frequency. Example at $X$:
\begin{center}
\begin{tabular}{|p{0.5\textwidth}|p{0.33\textwidth}|}
\hline
Interval & Root \\
\hline
below $1194.4\,\mathrm{Hz}$ & $835.6\,\mathrm{Hz}$ \\
exactly at $1194.4\,\mathrm{Hz}$ & sin pattern (node at the resonator) \\
between $1194.4\,\mathrm{Hz}$ and the next bare value & $1652.4\,\mathrm{Hz}$ \\
\hline
\end{tabular}
\end{center}

\paragraph{Consequence 2: the stuck modes.}
A $d$-fold bare value is one pole, so it supplies only one interval boundary. The
remaining $d-1$ combinations have $w_0 = 0$ and sit exactly on the pole. Counting
all roots gives all bare modes plus one extra root per $\mathbf{k}$: the extra
band of the one extra resonator degree of freedom.

\paragraph{Consequence 3: which side of $f_r$ a root is on.}
At a root, $Y = P$, and $Y > 0$ exactly when $\lambda < \lambda_r$. Therefore:
\begin{itemize}
  \item $P > 0$ at the root: the root is \emph{below} $f_r$; the plate at the
        attachment acts spring-like overall;
  \item $P < 0$ at the root: the root is \emph{above} $f_r$; the plate at the
        attachment acts mass-like overall.
\end{itemize}
Below the lowest bare value every term is positive, so $P > 0$ and band 1 lies
below $f_r$. At $\Gamma$, $P = -1 + S$, where $-1$ is the $\mathbf{G}=0$ wave
(rigid cell) and $S$ is the sum over all others. The band at $\Gamma$ sits above
$f_r$ only while $S < 1$.

% ---------------------------------------------------------------------
\section{The sum for Reissner--Mindlin}

In the code (\texttt{plane\_wave\_point\_compliance}) each wave enters through
$g(\kappa, \omega)$. For Kirchhoff,
\begin{equation}
  g = \frac{1}{D\kappa^4 - \rho h\,\omega^2} .
\end{equation}
For RM a wave has two coupled unknowns, the deflection and the rotation along
the wave, with dynamic stiffness matrix
\begin{equation}
  \begin{pmatrix}
    D_s\kappa^2 - \rho h\,\omega^2 & -D_s\kappa \\
    -D_s\kappa & D\kappa^2 + D_s - I_r\omega^2
  \end{pmatrix}
  =
  \begin{pmatrix} a & b \\ b & c \end{pmatrix} ,
\end{equation}
and $g$ is the deflection entry of its inverse, $g = c/(ac - b^2)$. At large
$\kappa$ the inertia terms are negligible:
\begin{align}
  ac - b^2 &= D_s\kappa^2\,(D\kappa^2 + D_s) - D_s^2\kappa^2 = D\,D_s\,\kappa^4 , \\
  g &\approx \frac{D\kappa^2 + D_s}{D\,D_s\,\kappa^4}
     = \frac{1}{D\kappa^4} + \frac{1}{D_s\kappa^2} .
\end{align}
Bending and shear compliance simply add. Kirchhoff is the limit
$D_s \to \infty$.

\paragraph{Why this matters for the sum.}
In 2D the number of waves with wavenumber between $\kappa$ and $\kappa +
\mathrm{d}\kappa$ is proportional to the area of the ring, $\kappa\,\mathrm{d}\kappa$.
The tail of the sum therefore behaves like
\begin{itemize}
  \item Kirchhoff: $\int \kappa\,\mathrm{d}\kappa/\kappa^4$, which \emph{converges};
  \item RM shear term: $\int \kappa\,\mathrm{d}\kappa/\kappa^2 = \int \mathrm{d}\kappa/\kappa = \ln\kappa$,
        which \emph{diverges slowly}.
\end{itemize}
Physically, a true point force on a shear-deformable plate gives infinite local
deflection. Real attachments have a finite size; the cutoff in the sum, or the
mesh in FE, plays that role.
\begin{center}
\begin{tabular}{|l|l|l|l|}
\hline
Cutoff & $\Gamma$ band 2, RM & $M$ band 1, RM & $\Gamma$ band 2, Kirchhoff \\
\hline
6 & 1276.1 & 1033.1 & 1280.2 \\
12 & 1274.7 & 1032.1 & 1279.8 \\
24 & 1273.5 & 1031.4 & 1279.7 \\
48 & 1272.4 & 1030.7 & 1279.6 \\
96 & 1271.3 & 1030.0 & -- \\
\hline
\end{tabular}
\end{center}
All values are in Hz; ``cutoff'' is the largest harmonic index kept in the sum.
Each doubling lowers the RM values by about the same amount
($\approx 1.1\,\mathrm{Hz}$ at $\Gamma$, $\approx 0.7\,\mathrm{Hz}$ at $M$), the
signature of a logarithm. The Kirchhoff values settle. FE shows the same effect
under mesh refinement ($1273.5$ / $1271.7$ / $1270.5\,\mathrm{Hz}$ for $10$/$20$/$40$
elements). It is about $0.1\,\%$ per doubling, so small in practice. The upper
gap edge at $X$ is not affected, because the sin pattern does not feel the point
force.

% ---------------------------------------------------------------------
\section{More resonators per cell}

Take resonators at positions $\mathbf{x}_j$, with forces $F_j$ and plate
deflections $w_j$. The steps are the same as for one resonator:
\begin{enumerate}
  \item Each wave feels all the forces:
        $(\lambda_n^0 - \lambda)\, q_n = \sum_l F_l\,\varphi_n(\mathbf{x}_l)^*$.
  \item Solve for $q_n$ and insert into
        $w_j = \sum_n \varphi_n(\mathbf{x}_j)\, q_n$:
        \begin{equation}
          w_j = \sum_l C_{jl}\, F_l ,
          \qquad
          C_{jl} = \sum_n \frac{\varphi_n(\mathbf{x}_j)\,\varphi_n(\mathbf{x}_l)^*}{\lambda_n^0 - \lambda} .
        \end{equation}
  \item Each resonator gives $F_l = Z_l\, w_l$. In matrix form $w = C Z w$, with
        $Z$ diagonal, so a nonzero solution needs
        \begin{equation}
          \boxed{\;\det(I - C Z) = 0\;}
        \end{equation}
\end{enumerate}
With one resonator this is $1 = ZC$. $C_{jl}$ is the deflection at resonator $j$
per unit force at resonator $l$. Each added resonator adds one band and can
couple to patterns that the others cannot.

\paragraph{Where the plane-wave model stops being valid.}
It assumes a uniform plate and point attachments. Cutouts, stiffeners or real
attachment patches need FE. Inside its validity it is nearly free to evaluate, so
it is both the check on FE and a fast design tool.

% ---------------------------------------------------------------------
\section{How detuning moves the edges}

\paragraph{Raising $f_r$.}
$Y = (1 - \lambda/\lambda_r)/\mu$ keeps its value $1/\mu$ at $\lambda = 0$, while
its zero moves to the right, so $Y$ increases at every $\lambda > 0$. At the old
root $Y - P$ is now positive, and since $Y - P$ falls with $\lambda$, the root
moves right: \emph{every coupled root rises with $f_r$}.

\paragraph{Raising $\mu$.}
$Y$ is zero at $\lambda = \lambda_r$ for any $\mu$, so raising $\mu$ turns the
line about that point. Below $\lambda_r$, $Y$ is positive and shrinks, so roots
below $f_r$ go \emph{down}. Above $\lambda_r$, $Y$ is negative and gets closer to
zero, so roots above $f_r$ go \emph{up}. A heavier resonator pushes roots
\emph{away from} $f_r$, which is level repulsion. Check at $f_r = 1200\,\mathrm{Hz}$,
$\mu = 0.25 \to 0.5$: band 3 at $X$ goes from $1653$ to $1842\,\mathrm{Hz}$, band 1
at $M$ from $1032$ to $927\,\mathrm{Hz}$.

\begin{center}
\begin{tabular}{|l|l|l|l|}
\hline
$f_r$ [Hz] & $M$ band 1 [Hz] & $\Gamma$ band 2 [Hz] & $\Gamma$ band 2 $/\, f_r$ \\
\hline
300 & 297 & 334 & 1.11 \\
800 & 748 & 874 & 1.09 \\
1200 & 1031 & 1273 & 1.06 \\
1500 & 1188 & 1546 & 1.03 \\
2000 & 1360 & 1940 & 0.97 \\
5000 & 1600 & 2957 & 0.59 \\
200000 & 1648 & 3279 & 0.02 \\
\hline
\end{tabular}
\end{center}

\paragraph{Limits.}
\begin{itemize}
  \item $f_r \to 0$: $Y$ falls very steeply, so the root sits at
        $\lambda \approx \lambda_r$, and band 1 at $M$ follows $f_r$.
  \item $f_r \to \infty$: $Y$ becomes the horizontal line $1/\mu$, so
        Eq.~\eqref{eq:Y_equals_P} is the rigid-mass condition. The spring is
        infinitely stiff, and band 1 at $M$ saturates at $1648\,\mathrm{Hz}$. The
        four-wave rigid-mass formula $2378/\sqrt{1+4\mu}$ gives
        $1682\,\mathrm{Hz}$; the other waves account for the rest.
\end{itemize}

\paragraph{Three tuning values mark the changes.}
\begin{center}
\begin{tabular}{|l|p{0.62\textwidth}|}
\hline
$f_r$ & What happens \\
\hline
$\approx 1118\,\mathrm{Hz}$ & $\Gamma$ band 2 reaches $f_X$, and the upper edge becomes pinned \\
$\approx 1513\,\mathrm{Hz}$ & $M$ band 1 reaches $f_X$, and the gap closes \\
$\approx 1765\,\mathrm{Hz}$ & $\Gamma$ band 2 equals $f_r$: $S$ reaches $1$, and the plate at $\Gamma$ turns from mass-like to spring-like \\
\hline
\end{tabular}
\end{center}
\paragraph{Why the gap is widest near $1118\,\mathrm{Hz}$.}
Below it, both edges move up with $f_r$, the upper one faster, so the width
grows. Above it, the upper edge is frozen while the lower one keeps rising, so
the width shrinks.

% ---------------------------------------------------------------------
\section{Two resonators per cell: numbers}
\label{sec:two_resonators}

Identical resonators, $f_r = 1200\,\mathrm{Hz}$, plane-wave model on
$\Gamma$--$X$--$M$--$Y$--$\Gamma$ with $15$ $\mathbf{k}$-points per segment. The
model reproduces the one-resonator values to about $1\,\mathrm{Hz}$, but it is not
FE. The path needs $Y$ because a resonator at $(a/2, 0)$ breaks the square
symmetry. The two-resonator gaps lie between bands 2 and 3, because each resonator
adds a band.
\begin{center}
\begin{tabular}{|p{0.23\textwidth}|p{0.15\textwidth}|p{0.18\textwidth}|l|l|}
\hline
Resonators per cell & At $X$ [Hz] & At $Y$ [Hz] & Gap [Hz] & Width \\
\hline
one at $(0,0)$, $\mu = 0.25$ & 835, \textbf{1194}, 1653 & 835, \textbf{1194}, 1653 & 1031--1194 & 163 \\
one at $(0,0)$, $\mu = 0.5$ & 724, \textbf{1194}, 1841 & 724, \textbf{1194}, 1841 & 926--1194 & 269 \\
$(0,0) + (a/2, 0)$, $\mu = 0.25$ each & 835 $\times 2$, 1653 $\times 2$ & 735, 1137, \textbf{1194}, 1913 & 1137--1194 & 58 \\
$(0,0) + (a/2, a/2)$, $\mu = 0.25$ each & 835 $\times 2$, 1653 $\times 2$ & 835 $\times 2$, 1653 $\times 2$ & 1114--1429 & 315 \\
\hline
\end{tabular}
\end{center}

\paragraph{The diagonal arrangement is a smaller lattice.}
Two identical resonators at $(0,0)$ and $(a/2, a/2)$ form a square lattice rotated
by $45^\circ$ with spacing $a/\sqrt{2}$. The first Bragg frequency moves up, from
$X$ at $1194\,\mathrm{Hz}$ to the old $M$ at $2378\,\mathrm{Hz}$. A second resonator
and a smaller cell are therefore the same physics here. The gap then contains
$f_r$, so it is a local-resonance gap.

\paragraph{Extra question: why is the $(0,0)+(a/2,0)$ gap so narrow?}
Sorted values at $Y$:
\begin{center}
\begin{tabular}{|l|l|p{0.58\textwidth}|}
\hline
Band at $Y$ & $f$ [Hz] & Mode \\
\hline
1 & 735 & $\cos(\pi y/a)$, both resonators moving together, lower branch \\
2 & \textbf{1137} & the two resonators moving in antiphase, new mode \\
3 & \textbf{1194.4} & $\sin(\pi y/a)$, node at both resonators, bare \\
4 & 1913 & $\cos(\pi y/a)$, both resonators moving together, upper branch \\
\hline
\end{tabular}
\end{center}
$1137\,\mathrm{Hz}$ is the \emph{maximum} of band 2, so it is the new lower edge;
the upper edge is the minimum of band 3, the pinned mode. The gap shrinks because
the second resonator adds a band, and at $Y$ that band lands inside the old gap.

Why it sits at $1137\,\mathrm{Hz}$: both resonators have $y = 0$, so the waves at
$Y$, which vary only in $y$, have the same value at both. If the resonators move
in antiphase, $F_1 = -F_2$, their forces on those waves cancel exactly, and the $Y$
waves stay at rest. The opposite forces do excite waves that alternate in $x$
between the two resonators: $(\pm 2\pi/a, \pm\pi/a)$, at about
$5868\,\mathrm{Hz}$, and higher waves with odd multiples of $2\pi/a$ in $x$. For
such a wave the two resonators contribute with opposite sign and the antiphase
forces add up, giving a factor $2$:
\begin{equation}
  1 = 2\,\mu_\mathrm{eff}\, P_\mathrm{odd},
  \qquad
  P_\mathrm{odd} = \sum_{\text{odd waves}} \frac{\lambda}{\lambda_n^0 - \lambda} .
\end{equation}
Insert $\mu_\mathrm{eff} = \mu/(1-\beta)$ and multiply by $(1-\beta)$:
\begin{equation}
  1 - \beta = 2\mu\, P_\mathrm{odd} .
\end{equation}
All of these waves lie far above $1137\,\mathrm{Hz}$, so $P_\mathrm{odd} > 0$ and
the plate acts like a soft spring, which pulls the mode below $f_r$. The four waves
at $5868\,\mathrm{Hz}$ give $P_\mathrm{odd} \approx 0.16$, and the higher waves
bring it to about $0.2$, so
\begin{equation}
  \beta \approx 1 - 2\cdot 0.25\cdot 0.205 = 0.897 ,
  \qquad
  f \approx 1200\cdot\sqrt{0.897} \approx 1137\,\mathrm{Hz}.
\end{equation}

% ---------------------------------------------------------------------
\section{Check questions with answers}

\paragraph{Question 1.} \emph{As $f_r$ goes up, band 1 at $M$ rises but saturates
near $1650\,\mathrm{Hz}$, well below its bare value of $2378\,\mathrm{Hz}$. Why does
it stop there, and what has the resonator physically become?}

\emph{Answer.} With $f_r \to \infty$ the spring cannot deform, so the resonator is
a rigid mass glued to the plate. Band 1 at $M$ settles at the mass-loaded value:
$1682\,\mathrm{Hz}$ from the four-wave formula, $1648\,\mathrm{Hz}$ with all waves.
A rigid mass can only lower a frequency, so stiffening the spring further changes
nothing.

\paragraph{Question 2.} \emph{A heavier resonator pushes roots away from $f_r$.
Where in the $Y$--$P$ picture do you see this, and what does it mean for a gap that
contains $f_r$ compared with one lying entirely below $f_r$?}

\emph{Answer.} Raising $\mu$ turns $Y$ about its zero at $\lambda = \lambda_r$.
Roots left of that point move left, roots right of it move right.
\begin{center}
\begin{tabular}{|p{0.25\textwidth}|l|l|p{0.17\textwidth}|}
\hline
Tuning & Gap at $\mu = 0.25$ & Gap at $\mu = 0.5$ & What moved \\
\hline
$f_r = 800\,\mathrm{Hz}$ (contains $f_r$) & 748--874\,Hz & 707--936\,Hz & both edges, in opposite directions \\
$f_r = 1200\,\mathrm{Hz}$ (below $f_r$) & 1032--1194\,Hz & 927--1194\,Hz & lower edge only \\
\hline
\end{tabular}
\end{center}
Widening does not make the $1200\,\mathrm{Hz}$ gap a local-resonance gap. A gap that
contains $f_r$ has one edge on each side of the pivot and is a true local-resonance
gap. The $1200\,\mathrm{Hz}$ gap lies entirely below it, and its upper edge is the
sin pattern at $X$, which is not a root of $Y = P$ at all because it has a node at
the resonator. It stays a Bragg gap with a resonance-boosted mass.

\paragraph{Question 3.} \emph{Add a second resonator in each cell at $(a/2, 0)$, same
$f_r$. What happens to the sin pattern at $X$, and would you expect the upper edge to
stay pinned at $1194\,\mathrm{Hz}$?}

\emph{Answer.} At $X$, $\cos(\pi x/a)$ has its antinode at the first resonator and a
node at the second, and $\sin(\pi x/a)$ the reverse. That makes two identical
plate-plus-resonator pairs. The sin pattern does not simply rise; it \emph{splits}
into a lower and an upper branch, like cos, and the $1194.4\,\mathrm{Hz}$ mode is
gone from $X$.

But $X$ is not the only point that matters. At $Y$ the pattern $\sin(\pi y/a)$ is
zero at $y = 0$, where both resonators sit, so it stays at $1194.4\,\mathrm{Hz}$ and
the upper edge stays pinned. Moving the second resonator to the cell centre removes
the pinning at both $X$ and $Y$ (Section~\ref{sec:two_resonators}).

% =====================================================================

\chapter{The plane-wave condition as a general tool}
\label{ch:pwe_general}

This chapter takes the plane-wave expansion (PWE) condition of
Chapter~\ref{ch:gap_position} and asks what can be read off it without solving
it: how many bands exist, where they can lie, what happens to them when the
resonator is detuned, and what changes with more than one resonator per cell.
Numbers come from the Reissner--Mindlin (RM) plane-wave model with $\mu = 0.25$
and $f_r = 1200\,\mathrm{Hz}$ unless stated otherwise.

\begin{center}
\begin{tabular}{|l|p{0.62\textwidth}|}
\hline
Symbol & Meaning \\
\hline
$\lambda = \omega^2$ & eigenvalue of one dot in the diagram \\
$\lambda_n^0$ & bare plate eigenvalue of wave $n$ at this $\mathbf{k}$ \\
$\lambda_r = k_s/m$ & the resonator's own eigenvalue ($k_s$ spring, $m$ mass) \\
$\mu = m/(\rho h A)$ & resonator mass over plate mass per cell ($\rho h$ mass per area, $A$ cell area) \\
$\beta = \lambda/\lambda_r$ & squared frequency relative to the resonator \\
$w_0$ & plate deflection at the resonator \\
$Z$, $C$ & force per deflection of the resonator; deflection per force of the plate \\
$P(\lambda)$ & plate curve, Eq.~\eqref{eq:plate_curve} \\
$Y(\lambda)$ & resonator line, Eq.~\eqref{eq:resonator_line} \\
$S$ & at $\Gamma$: the part of $P$ from all waves except $\mathbf{G}=0$ \\
$g(\kappa, \omega)$ & deflection per unit point force of one plane wave of wavenumber $\kappa$ \\
$D$, $D_s$ & bending stiffness and shear stiffness of the plate \\
$I_r = \rho h^3/12$ & rotary inertia per area \\
$f_X$ & bare plate frequency at $X$, $1194.4\,\mathrm{Hz}$ \\
\hline
\end{tabular}
\end{center}

% ---------------------------------------------------------------------
\section{Recap of the condition}

Three facts combine:
\begin{itemize}
  \item A resonator exerts a force on the plate like a mass that depends on
        frequency: $F = m_\mathrm{eff}\,\lambda\, w_0$, with
        $m_\mathrm{eff} = m/(1-\beta)$.
  \item A point force excites every mass-normalized bare wave $\varphi_n$ with
        amplitude
        \[ q_n = \frac{F\,\varphi_n(\mathbf{x}_0)^*}{\lambda_n^0 - \lambda} . \]
  \item Adding up all waves at $\mathbf{x}_0$ must give back $w_0$.
\end{itemize}
Inserting one into the other, with
$|\varphi_n(\mathbf{x}_0)|^2 = 1/(\rho h A)$ and
$\mu_\mathrm{eff} = \mu/(1-\beta)$, gave
\begin{equation}
  1 = \mu_\mathrm{eff}\,\sum_n \frac{\lambda}{\lambda_n^0 - \lambda} .
  \label{eq:pwe_recap}
\end{equation}

% ---------------------------------------------------------------------
\section{Separating resonator and plate}

Name the sum:
\begin{equation}
  P(\lambda) = \sum_n \frac{\lambda}{\lambda_n^0 - \lambda} .
  \label{eq:plate_curve}
\end{equation}
Then Eq.~\eqref{eq:pwe_recap} reads $1 = \mu P/(1-\beta)$. Multiply both sides
by $(1-\beta)/\mu$:
\begin{equation}
  \frac{1-\beta}{\mu} = P(\lambda) .
\end{equation}
Name the left side:
\begin{equation}
  Y(\lambda) = \frac{1 - \lambda/\lambda_r}{\mu} .
  \label{eq:resonator_line}
\end{equation}
The condition is then simply
\begin{equation}
  \boxed{\;Y(\lambda) = P(\lambda)\;}
  \label{eq:Y_equals_P}
\end{equation}
The two sides have separate origins:
\begin{itemize}
  \item $Y$ contains only the resonator ($\lambda_r$, $\mu$). It is a falling
        straight line in $\lambda$, with $Y = 1/\mu$ at $\lambda = 0$ and
        $Y = 0$ at $\lambda = \lambda_r$.
  \item $P$ contains only the plate and $\mathbf{k}$. It does not know about
        the resonator.
\end{itemize}

\paragraph{$P$ rises everywhere.} With the quotient rule, for one term,
\begin{equation}
  \frac{\mathrm{d}}{\mathrm{d}\lambda}\,\frac{\lambda}{\lambda_n^0 - \lambda}
  = \frac{(\lambda_n^0 - \lambda) + \lambda}{(\lambda_n^0 - \lambda)^2}
  = \frac{\lambda_n^0}{(\lambda_n^0 - \lambda)^2} \ge 0 .
\end{equation}

\paragraph{$P$ has a pole at every distinct bare value.}
Just below $\lambda_n^0$ the denominator is small and positive, so
$P \to +\infty$. Just above it, $P \to -\infty$. Between two neighbouring bare
values, $P$ therefore climbs from $-\infty$ to $+\infty$.

The dots of the diagram at this $\mathbf{k}$ are the intersections of a falling
line with these rising branches.

% ---------------------------------------------------------------------
\section{What follows from that picture}

Consider $Y - P$ between two neighbouring bare values. It is a falling function
minus a rising one, so it falls strictly, from $+\infty$ to $-\infty$, and
crosses zero \emph{exactly once}. The same holds below the lowest bare value,
where $Y - P$ starts positive.

\paragraph{Consequence 1: bands are trapped.}
There is exactly one root between neighbouring bare frequencies. A single point
resonator cannot push a band past a bare frequency. Example at $X$:
\begin{center}
\begin{tabular}{|p{0.5\textwidth}|p{0.33\textwidth}|}
\hline
Interval & Root \\
\hline
below $1194.4\,\mathrm{Hz}$ & $835.6\,\mathrm{Hz}$ \\
exactly at $1194.4\,\mathrm{Hz}$ & sin pattern (node at the resonator) \\
between $1194.4\,\mathrm{Hz}$ and the next bare value & $1652.4\,\mathrm{Hz}$ \\
\hline
\end{tabular}
\end{center}

\paragraph{Consequence 2: the stuck modes.}
A $d$-fold bare value is one pole, so it supplies only one interval boundary. The
remaining $d-1$ combinations have $w_0 = 0$ and sit exactly on the pole. Counting
all roots gives all bare modes plus one extra root per $\mathbf{k}$: the extra
band of the one extra resonator degree of freedom.

\paragraph{Consequence 3: which side of $f_r$ a root is on.}
At a root, $Y = P$, and $Y > 0$ exactly when $\lambda < \lambda_r$. Therefore:
\begin{itemize}
  \item $P > 0$ at the root: the root is \emph{below} $f_r$; the plate at the
        attachment acts spring-like overall;
  \item $P < 0$ at the root: the root is \emph{above} $f_r$; the plate at the
        attachment acts mass-like overall.
\end{itemize}
Below the lowest bare value every term is positive, so $P > 0$ and band 1 lies
below $f_r$. At $\Gamma$, $P = -1 + S$, where $-1$ is the $\mathbf{G}=0$ wave
(rigid cell) and $S$ is the sum over all others. The band at $\Gamma$ sits above
$f_r$ only while $S < 1$.

% ---------------------------------------------------------------------
\section{The sum for Reissner--Mindlin}

In the code (\texttt{plane\_wave\_point\_compliance}) each wave enters through
$g(\kappa, \omega)$. For Kirchhoff,
\begin{equation}
  g = \frac{1}{D\kappa^4 - \rho h\,\omega^2} .
\end{equation}
For RM a wave has two coupled unknowns, the deflection and the rotation along
the wave, with dynamic stiffness matrix
\begin{equation}
  \begin{pmatrix}
    D_s\kappa^2 - \rho h\,\omega^2 & -D_s\kappa \\
    -D_s\kappa & D\kappa^2 + D_s - I_r\omega^2
  \end{pmatrix}
  =
  \begin{pmatrix} a & b \\ b & c \end{pmatrix} ,
\end{equation}
and $g$ is the deflection entry of its inverse, $g = c/(ac - b^2)$. At large
$\kappa$ the inertia terms are negligible:
\begin{align}
  ac - b^2 &= D_s\kappa^2\,(D\kappa^2 + D_s) - D_s^2\kappa^2 = D\,D_s\,\kappa^4 , \\
  g &\approx \frac{D\kappa^2 + D_s}{D\,D_s\,\kappa^4}
     = \frac{1}{D\kappa^4} + \frac{1}{D_s\kappa^2} .
\end{align}
Bending and shear compliance simply add. Kirchhoff is the limit
$D_s \to \infty$.

\paragraph{Why this matters for the sum.}
In 2D the number of waves with wavenumber between $\kappa$ and $\kappa +
\mathrm{d}\kappa$ is proportional to the area of the ring, $\kappa\,\mathrm{d}\kappa$.
The tail of the sum therefore behaves like
\begin{itemize}
  \item Kirchhoff: $\int \kappa\,\mathrm{d}\kappa/\kappa^4$, which \emph{converges};
  \item RM shear term: $\int \kappa\,\mathrm{d}\kappa/\kappa^2 = \int \mathrm{d}\kappa/\kappa = \ln\kappa$,
        which \emph{diverges slowly}.
\end{itemize}
Physically, a true point force on a shear-deformable plate gives infinite local
deflection. Real attachments have a finite size; the cutoff in the sum, or the
mesh in FE, plays that role.
\begin{center}
\begin{tabular}{|l|l|l|l|}
\hline
Cutoff & $\Gamma$ band 2, RM & $M$ band 1, RM & $\Gamma$ band 2, Kirchhoff \\
\hline
6 & 1276.1 & 1033.1 & 1280.2 \\
12 & 1274.7 & 1032.1 & 1279.8 \\
24 & 1273.5 & 1031.4 & 1279.7 \\
48 & 1272.4 & 1030.7 & 1279.6 \\
96 & 1271.3 & 1030.0 & -- \\
\hline
\end{tabular}
\end{center}
All values are in Hz; ``cutoff'' is the largest harmonic index kept in the sum.
Each doubling lowers the RM values by about the same amount
($\approx 1.1\,\mathrm{Hz}$ at $\Gamma$, $\approx 0.7\,\mathrm{Hz}$ at $M$), the
signature of a logarithm. The Kirchhoff values settle. FE shows the same effect
under mesh refinement ($1273.5$ / $1271.7$ / $1270.5\,\mathrm{Hz}$ for $10$/$20$/$40$
elements). It is about $0.1\,\%$ per doubling, so small in practice. The upper
gap edge at $X$ is not affected, because the sin pattern does not feel the point
force.

% ---------------------------------------------------------------------
\section{More resonators per cell}

Take resonators at positions $\mathbf{x}_j$, with forces $F_j$ and plate
deflections $w_j$. The steps are the same as for one resonator:
\begin{enumerate}
  \item Each wave feels all the forces:
        $(\lambda_n^0 - \lambda)\, q_n = \sum_l F_l\,\varphi_n(\mathbf{x}_l)^*$.
  \item Solve for $q_n$ and insert into
        $w_j = \sum_n \varphi_n(\mathbf{x}_j)\, q_n$:
        \begin{equation}
          w_j = \sum_l C_{jl}\, F_l ,
          \qquad
          C_{jl} = \sum_n \frac{\varphi_n(\mathbf{x}_j)\,\varphi_n(\mathbf{x}_l)^*}{\lambda_n^0 - \lambda} .
        \end{equation}
  \item Each resonator gives $F_l = Z_l\, w_l$. In matrix form $w = C Z w$, with
        $Z$ diagonal, so a nonzero solution needs
        \begin{equation}
          \boxed{\;\det(I - C Z) = 0\;}
        \end{equation}
\end{enumerate}
With one resonator this is $1 = ZC$. $C_{jl}$ is the deflection at resonator $j$
per unit force at resonator $l$. Each added resonator adds one band and can
couple to patterns that the others cannot.

\paragraph{Where the plane-wave model stops being valid.}
It assumes a uniform plate and point attachments. Cutouts, stiffeners or real
attachment patches need FE. Inside its validity it is nearly free to evaluate, so
it is both the check on FE and a fast design tool.

% ---------------------------------------------------------------------
\section{How detuning moves the edges}

\paragraph{Raising $f_r$.}
$Y = (1 - \lambda/\lambda_r)/\mu$ keeps its value $1/\mu$ at $\lambda = 0$, while
its zero moves to the right, so $Y$ increases at every $\lambda > 0$. At the old
root $Y - P$ is now positive, and since $Y - P$ falls with $\lambda$, the root
moves right: \emph{every coupled root rises with $f_r$}.

\paragraph{Raising $\mu$.}
$Y$ is zero at $\lambda = \lambda_r$ for any $\mu$, so raising $\mu$ turns the
line about that point. Below $\lambda_r$, $Y$ is positive and shrinks, so roots
below $f_r$ go \emph{down}. Above $\lambda_r$, $Y$ is negative and gets closer to
zero, so roots above $f_r$ go \emph{up}. A heavier resonator pushes roots
\emph{away from} $f_r$, which is level repulsion. Check at $f_r = 1200\,\mathrm{Hz}$,
$\mu = 0.25 \to 0.5$: band 3 at $X$ goes from $1653$ to $1842\,\mathrm{Hz}$, band 1
at $M$ from $1032$ to $927\,\mathrm{Hz}$.

\begin{center}
\begin{tabular}{|l|l|l|l|}
\hline
$f_r$ [Hz] & $M$ band 1 [Hz] & $\Gamma$ band 2 [Hz] & $\Gamma$ band 2 $/\, f_r$ \\
\hline
300 & 297 & 334 & 1.11 \\
800 & 748 & 874 & 1.09 \\
1200 & 1031 & 1273 & 1.06 \\
1500 & 1188 & 1546 & 1.03 \\
2000 & 1360 & 1940 & 0.97 \\
5000 & 1600 & 2957 & 0.59 \\
200000 & 1648 & 3279 & 0.02 \\
\hline
\end{tabular}
\end{center}

\paragraph{Limits.}
\begin{itemize}
  \item $f_r \to 0$: $Y$ falls very steeply, so the root sits at
        $\lambda \approx \lambda_r$, and band 1 at $M$ follows $f_r$.
  \item $f_r \to \infty$: $Y$ becomes the horizontal line $1/\mu$, so
        Eq.~\eqref{eq:Y_equals_P} is the rigid-mass condition. The spring is
        infinitely stiff, and band 1 at $M$ saturates at $1648\,\mathrm{Hz}$. The
        four-wave rigid-mass formula $2378/\sqrt{1+4\mu}$ gives
        $1682\,\mathrm{Hz}$; the other waves account for the rest.
\end{itemize}

\paragraph{Three tuning values mark the changes.}
\begin{center}
\begin{tabular}{|l|p{0.62\textwidth}|}
\hline
$f_r$ & What happens \\
\hline
$\approx 1118\,\mathrm{Hz}$ & $\Gamma$ band 2 reaches $f_X$, and the upper edge becomes pinned \\
$\approx 1513\,\mathrm{Hz}$ & $M$ band 1 reaches $f_X$, and the gap closes \\
$\approx 1765\,\mathrm{Hz}$ & $\Gamma$ band 2 equals $f_r$: $S$ reaches $1$, and the plate at $\Gamma$ turns from mass-like to spring-like \\
\hline
\end{tabular}
\end{center}
\paragraph{Why the gap is widest near $1118\,\mathrm{Hz}$.}
Below it, both edges move up with $f_r$, the upper one faster, so the width
grows. Above it, the upper edge is frozen while the lower one keeps rising, so
the width shrinks.

% ---------------------------------------------------------------------
\section{Two resonators per cell: numbers}
\label{sec:two_resonators}

Identical resonators, $f_r = 1200\,\mathrm{Hz}$, plane-wave model on
$\Gamma$--$X$--$M$--$Y$--$\Gamma$ with $15$ $\mathbf{k}$-points per segment. The
model reproduces the one-resonator values to about $1\,\mathrm{Hz}$, but it is not
FE. The path needs $Y$ because a resonator at $(a/2, 0)$ breaks the square
symmetry. The two-resonator gaps lie between bands 2 and 3, because each resonator
adds a band.
\begin{center}
\begin{tabular}{|p{0.23\textwidth}|p{0.15\textwidth}|p{0.18\textwidth}|l|l|}
\hline
Resonators per cell & At $X$ [Hz] & At $Y$ [Hz] & Gap [Hz] & Width \\
\hline
one at $(0,0)$, $\mu = 0.25$ & 835, \textbf{1194}, 1653 & 835, \textbf{1194}, 1653 & 1031--1194 & 163 \\
one at $(0,0)$, $\mu = 0.5$ & 724, \textbf{1194}, 1841 & 724, \textbf{1194}, 1841 & 926--1194 & 269 \\
$(0,0) + (a/2, 0)$, $\mu = 0.25$ each & 835 $\times 2$, 1653 $\times 2$ & 735, 1137, \textbf{1194}, 1913 & 1137--1194 & 58 \\
$(0,0) + (a/2, a/2)$, $\mu = 0.25$ each & 835 $\times 2$, 1653 $\times 2$ & 835 $\times 2$, 1653 $\times 2$ & 1114--1429 & 315 \\
\hline
\end{tabular}
\end{center}

\paragraph{The diagonal arrangement is a smaller lattice.}
Two identical resonators at $(0,0)$ and $(a/2, a/2)$ form a square lattice rotated
by $45^\circ$ with spacing $a/\sqrt{2}$. The first Bragg frequency moves up, from
$X$ at $1194\,\mathrm{Hz}$ to the old $M$ at $2378\,\mathrm{Hz}$. A second resonator
and a smaller cell are therefore the same physics here. The gap then contains
$f_r$, so it is a local-resonance gap.

\paragraph{Extra question: why is the $(0,0)+(a/2,0)$ gap so narrow?}
Sorted values at $Y$:
\begin{center}
\begin{tabular}{|l|l|p{0.58\textwidth}|}
\hline
Band at $Y$ & $f$ [Hz] & Mode \\
\hline
1 & 735 & $\cos(\pi y/a)$, both resonators moving together, lower branch \\
2 & \textbf{1137} & the two resonators moving in antiphase, new mode \\
3 & \textbf{1194.4} & $\sin(\pi y/a)$, node at both resonators, bare \\
4 & 1913 & $\cos(\pi y/a)$, both resonators moving together, upper branch \\
\hline
\end{tabular}
\end{center}
$1137\,\mathrm{Hz}$ is the \emph{maximum} of band 2, so it is the new lower edge;
the upper edge is the minimum of band 3, the pinned mode. The gap shrinks because
the second resonator adds a band, and at $Y$ that band lands inside the old gap.

Why it sits at $1137\,\mathrm{Hz}$: both resonators have $y = 0$, so the waves at
$Y$, which vary only in $y$, have the same value at both. If the resonators move
in antiphase, $F_1 = -F_2$, their forces on those waves cancel exactly, and the $Y$
waves stay at rest. The opposite forces do excite waves that alternate in $x$
between the two resonators: $(\pm 2\pi/a, \pm\pi/a)$, at about
$5868\,\mathrm{Hz}$, and higher waves with odd multiples of $2\pi/a$ in $x$. For
such a wave the two resonators contribute with opposite sign and the antiphase
forces add up, giving a factor $2$:
\begin{equation}
  1 = 2\,\mu_\mathrm{eff}\, P_\mathrm{odd},
  \qquad
  P_\mathrm{odd} = \sum_{\text{odd waves}} \frac{\lambda}{\lambda_n^0 - \lambda} .
\end{equation}
Insert $\mu_\mathrm{eff} = \mu/(1-\beta)$ and multiply by $(1-\beta)$:
\begin{equation}
  1 - \beta = 2\mu\, P_\mathrm{odd} .
\end{equation}
All of these waves lie far above $1137\,\mathrm{Hz}$, so $P_\mathrm{odd} > 0$ and
the plate acts like a soft spring, which pulls the mode below $f_r$. The four waves
at $5868\,\mathrm{Hz}$ give $P_\mathrm{odd} \approx 0.16$, and the higher waves
bring it to about $0.2$, so
\begin{equation}
  \beta \approx 1 - 2\cdot 0.25\cdot 0.205 = 0.897 ,
  \qquad
  f \approx 1200\cdot\sqrt{0.897} \approx 1137\,\mathrm{Hz}.
\end{equation}

% ---------------------------------------------------------------------
\section{Check questions with answers}

\paragraph{Question 1.} \emph{As $f_r$ goes up, band 1 at $M$ rises but saturates
near $1650\,\mathrm{Hz}$, well below its bare value of $2378\,\mathrm{Hz}$. Why does
it stop there, and what has the resonator physically become?}

\emph{Answer.} With $f_r \to \infty$ the spring cannot deform, so the resonator is
a rigid mass glued to the plate. Band 1 at $M$ settles at the mass-loaded value:
$1682\,\mathrm{Hz}$ from the four-wave formula, $1648\,\mathrm{Hz}$ with all waves.
A rigid mass can only lower a frequency, so stiffening the spring further changes
nothing.

\paragraph{Question 2.} \emph{A heavier resonator pushes roots away from $f_r$.
Where in the $Y$--$P$ picture do you see this, and what does it mean for a gap that
contains $f_r$ compared with one lying entirely below $f_r$?}

\emph{Answer.} Raising $\mu$ turns $Y$ about its zero at $\lambda = \lambda_r$.
Roots left of that point move left, roots right of it move right.
\begin{center}
\begin{tabular}{|p{0.25\textwidth}|l|l|p{0.17\textwidth}|}
\hline
Tuning & Gap at $\mu = 0.25$ & Gap at $\mu = 0.5$ & What moved \\
\hline
$f_r = 800\,\mathrm{Hz}$ (contains $f_r$) & 748--874\,Hz & 707--936\,Hz & both edges, in opposite directions \\
$f_r = 1200\,\mathrm{Hz}$ (below $f_r$) & 1032--1194\,Hz & 927--1194\,Hz & lower edge only \\
\hline
\end{tabular}
\end{center}
Widening does not make the $1200\,\mathrm{Hz}$ gap a local-resonance gap. A gap that
contains $f_r$ has one edge on each side of the pivot and is a true local-resonance
gap. The $1200\,\mathrm{Hz}$ gap lies entirely below it, and its upper edge is the
sin pattern at $X$, which is not a root of $Y = P$ at all because it has a node at
the resonator. It stays a Bragg gap with a resonance-boosted mass.

\paragraph{Question 3.} \emph{Add a second resonator in each cell at $(a/2, 0)$, same
$f_r$. What happens to the sin pattern at $X$, and would you expect the upper edge to
stay pinned at $1194\,\mathrm{Hz}$?}

\emph{Answer.} At $X$, $\cos(\pi x/a)$ has its antinode at the first resonator and a
node at the second, and $\sin(\pi x/a)$ the reverse. That makes two identical
plate-plus-resonator pairs. The sin pattern does not simply rise; it \emph{splits}
into a lower and an upper branch, like cos, and the $1194.4\,\mathrm{Hz}$ mode is
gone from $X$.

But $X$ is not the only point that matters. At $Y$ the pattern $\sin(\pi y/a)$ is
zero at $y = 0$, where both resonators sit, so it stays at $1194.4\,\mathrm{Hz}$ and
the upper edge stays pinned. Moving the second resonator to the cell centre removes
the pinning at both $X$ and $Y$ (Section~\ref{sec:two_resonators}).

% =====================================================================

\chapter{The plane-wave condition as a general tool}
\label{ch:pwe_general}

This chapter takes the plane-wave expansion (PWE) condition of
Chapter~\ref{ch:gap_position} and asks what can be read off it without solving
it: how many bands exist, where they can lie, what happens to them when the
resonator is detuned, and what changes with more than one resonator per cell.
Numbers come from the Reissner--Mindlin (RM) plane-wave model with $\mu = 0.25$
and $f_r = 1200\,\mathrm{Hz}$ unless stated otherwise.

\begin{center}
\begin{tabular}{|l|p{0.62\textwidth}|}
\hline
Symbol & Meaning \\
\hline
$\lambda = \omega^2$ & eigenvalue of one dot in the diagram \\
$\lambda_n^0$ & bare plate eigenvalue of wave $n$ at this $\mathbf{k}$ \\
$\lambda_r = k_s/m$ & the resonator's own eigenvalue ($k_s$ spring, $m$ mass) \\
$\mu = m/(\rho h A)$ & resonator mass over plate mass per cell ($\rho h$ mass per area, $A$ cell area) \\
$\beta = \lambda/\lambda_r$ & squared frequency relative to the resonator \\
$w_0$ & plate deflection at the resonator \\
$Z$, $C$ & force per deflection of the resonator; deflection per force of the plate \\
$P(\lambda)$ & plate curve, Eq.~\eqref{eq:plate_curve} \\
$Y(\lambda)$ & resonator line, Eq.~\eqref{eq:resonator_line} \\
$S$ & at $\Gamma$: the part of $P$ from all waves except $\mathbf{G}=0$ \\
$g(\kappa, \omega)$ & deflection per unit point force of one plane wave of wavenumber $\kappa$ \\
$D$, $D_s$ & bending stiffness and shear stiffness of the plate \\
$I_r = \rho h^3/12$ & rotary inertia per area \\
$f_X$ & bare plate frequency at $X$, $1194.4\,\mathrm{Hz}$ \\
\hline
\end{tabular}
\end{center}

% ---------------------------------------------------------------------
\section{Recap of the condition}

Three facts combine:
\begin{itemize}
  \item A resonator exerts a force on the plate like a mass that depends on
        frequency: $F = m_\mathrm{eff}\,\lambda\, w_0$, with
        $m_\mathrm{eff} = m/(1-\beta)$.
  \item A point force excites every mass-normalized bare wave $\varphi_n$ with
        amplitude
        \[ q_n = \frac{F\,\varphi_n(\mathbf{x}_0)^*}{\lambda_n^0 - \lambda} . \]
  \item Adding up all waves at $\mathbf{x}_0$ must give back $w_0$.
\end{itemize}
Inserting one into the other, with
$|\varphi_n(\mathbf{x}_0)|^2 = 1/(\rho h A)$ and
$\mu_\mathrm{eff} = \mu/(1-\beta)$, gave
\begin{equation}
  1 = \mu_\mathrm{eff}\,\sum_n \frac{\lambda}{\lambda_n^0 - \lambda} .
  \label{eq:pwe_recap}
\end{equation}

% ---------------------------------------------------------------------
\section{Separating resonator and plate}

Name the sum:
\begin{equation}
  P(\lambda) = \sum_n \frac{\lambda}{\lambda_n^0 - \lambda} .
  \label{eq:plate_curve}
\end{equation}
Then Eq.~\eqref{eq:pwe_recap} reads $1 = \mu P/(1-\beta)$. Multiply both sides
by $(1-\beta)/\mu$:
\begin{equation}
  \frac{1-\beta}{\mu} = P(\lambda) .
\end{equation}
Name the left side:
\begin{equation}
  Y(\lambda) = \frac{1 - \lambda/\lambda_r}{\mu} .
  \label{eq:resonator_line}
\end{equation}
The condition is then simply
\begin{equation}
  \boxed{\;Y(\lambda) = P(\lambda)\;}
  \label{eq:Y_equals_P}
\end{equation}
The two sides have separate origins:
\begin{itemize}
  \item $Y$ contains only the resonator ($\lambda_r$, $\mu$). It is a falling
        straight line in $\lambda$, with $Y = 1/\mu$ at $\lambda = 0$ and
        $Y = 0$ at $\lambda = \lambda_r$.
  \item $P$ contains only the plate and $\mathbf{k}$. It does not know about
        the resonator.
\end{itemize}

\paragraph{$P$ rises everywhere.} With the quotient rule, for one term,
\begin{equation}
  \frac{\mathrm{d}}{\mathrm{d}\lambda}\,\frac{\lambda}{\lambda_n^0 - \lambda}
  = \frac{(\lambda_n^0 - \lambda) + \lambda}{(\lambda_n^0 - \lambda)^2}
  = \frac{\lambda_n^0}{(\lambda_n^0 - \lambda)^2} \ge 0 .
\end{equation}

\paragraph{$P$ has a pole at every distinct bare value.}
Just below $\lambda_n^0$ the denominator is small and positive, so
$P \to +\infty$. Just above it, $P \to -\infty$. Between two neighbouring bare
values, $P$ therefore climbs from $-\infty$ to $+\infty$.

The dots of the diagram at this $\mathbf{k}$ are the intersections of a falling
line with these rising branches.

% ---------------------------------------------------------------------
\section{What follows from that picture}

Consider $Y - P$ between two neighbouring bare values. It is a falling function
minus a rising one, so it falls strictly, from $+\infty$ to $-\infty$, and
crosses zero \emph{exactly once}. The same holds below the lowest bare value,
where $Y - P$ starts positive.

\paragraph{Consequence 1: bands are trapped.}
There is exactly one root between neighbouring bare frequencies. A single point
resonator cannot push a band past a bare frequency. Example at $X$:
\begin{center}
\begin{tabular}{|p{0.5\textwidth}|p{0.33\textwidth}|}
\hline
Interval & Root \\
\hline
below $1194.4\,\mathrm{Hz}$ & $835.6\,\mathrm{Hz}$ \\
exactly at $1194.4\,\mathrm{Hz}$ & sin pattern (node at the resonator) \\
between $1194.4\,\mathrm{Hz}$ and the next bare value & $1652.4\,\mathrm{Hz}$ \\
\hline
\end{tabular}
\end{center}

\paragraph{Consequence 2: the stuck modes.}
A $d$-fold bare value is one pole, so it supplies only one interval boundary. The
remaining $d-1$ combinations have $w_0 = 0$ and sit exactly on the pole. Counting
all roots gives all bare modes plus one extra root per $\mathbf{k}$: the extra
band of the one extra resonator degree of freedom.

\paragraph{Consequence 3: which side of $f_r$ a root is on.}
At a root, $Y = P$, and $Y > 0$ exactly when $\lambda < \lambda_r$. Therefore:
\begin{itemize}
  \item $P > 0$ at the root: the root is \emph{below} $f_r$; the plate at the
        attachment acts spring-like overall;
  \item $P < 0$ at the root: the root is \emph{above} $f_r$; the plate at the
        attachment acts mass-like overall.
\end{itemize}
Below the lowest bare value every term is positive, so $P > 0$ and band 1 lies
below $f_r$. At $\Gamma$, $P = -1 + S$, where $-1$ is the $\mathbf{G}=0$ wave
(rigid cell) and $S$ is the sum over all others. The band at $\Gamma$ sits above
$f_r$ only while $S < 1$.

% ---------------------------------------------------------------------
\section{The sum for Reissner--Mindlin}

In the code (\texttt{plane\_wave\_point\_compliance}) each wave enters through
$g(\kappa, \omega)$. For Kirchhoff,
\begin{equation}
  g = \frac{1}{D\kappa^4 - \rho h\,\omega^2} .
\end{equation}
For RM a wave has two coupled unknowns, the deflection and the rotation along
the wave, with dynamic stiffness matrix
\begin{equation}
  \begin{pmatrix}
    D_s\kappa^2 - \rho h\,\omega^2 & -D_s\kappa \\
    -D_s\kappa & D\kappa^2 + D_s - I_r\omega^2
  \end{pmatrix}
  =
  \begin{pmatrix} a & b \\ b & c \end{pmatrix} ,
\end{equation}
and $g$ is the deflection entry of its inverse, $g = c/(ac - b^2)$. At large
$\kappa$ the inertia terms are negligible:
\begin{align}
  ac - b^2 &= D_s\kappa^2\,(D\kappa^2 + D_s) - D_s^2\kappa^2 = D\,D_s\,\kappa^4 , \\
  g &\approx \frac{D\kappa^2 + D_s}{D\,D_s\,\kappa^4}
     = \frac{1}{D\kappa^4} + \frac{1}{D_s\kappa^2} .
\end{align}
Bending and shear compliance simply add. Kirchhoff is the limit
$D_s \to \infty$.

\paragraph{Why this matters for the sum.}
In 2D the number of waves with wavenumber between $\kappa$ and $\kappa +
\mathrm{d}\kappa$ is proportional to the area of the ring, $\kappa\,\mathrm{d}\kappa$.
The tail of the sum therefore behaves like
\begin{itemize}
  \item Kirchhoff: $\int \kappa\,\mathrm{d}\kappa/\kappa^4$, which \emph{converges};
  \item RM shear term: $\int \kappa\,\mathrm{d}\kappa/\kappa^2 = \int \mathrm{d}\kappa/\kappa = \ln\kappa$,
        which \emph{diverges slowly}.
\end{itemize}
Physically, a true point force on a shear-deformable plate gives infinite local
deflection. Real attachments have a finite size; the cutoff in the sum, or the
mesh in FE, plays that role.
\begin{center}
\begin{tabular}{|l|l|l|l|}
\hline
Cutoff & $\Gamma$ band 2, RM & $M$ band 1, RM & $\Gamma$ band 2, Kirchhoff \\
\hline
6 & 1276.1 & 1033.1 & 1280.2 \\
12 & 1274.7 & 1032.1 & 1279.8 \\
24 & 1273.5 & 1031.4 & 1279.7 \\
48 & 1272.4 & 1030.7 & 1279.6 \\
96 & 1271.3 & 1030.0 & -- \\
\hline
\end{tabular}
\end{center}
All values are in Hz; ``cutoff'' is the largest harmonic index kept in the sum.
Each doubling lowers the RM values by about the same amount
($\approx 1.1\,\mathrm{Hz}$ at $\Gamma$, $\approx 0.7\,\mathrm{Hz}$ at $M$), the
signature of a logarithm. The Kirchhoff values settle. FE shows the same effect
under mesh refinement ($1273.5$ / $1271.7$ / $1270.5\,\mathrm{Hz}$ for $10$/$20$/$40$
elements). It is about $0.1\,\%$ per doubling, so small in practice. The upper
gap edge at $X$ is not affected, because the sin pattern does not feel the point
force.

% ---------------------------------------------------------------------
\section{More resonators per cell}

Take resonators at positions $\mathbf{x}_j$, with forces $F_j$ and plate
deflections $w_j$. The steps are the same as for one resonator:
\begin{enumerate}
  \item Each wave feels all the forces:
        $(\lambda_n^0 - \lambda)\, q_n = \sum_l F_l\,\varphi_n(\mathbf{x}_l)^*$.
  \item Solve for $q_n$ and insert into
        $w_j = \sum_n \varphi_n(\mathbf{x}_j)\, q_n$:
        \begin{equation}
          w_j = \sum_l C_{jl}\, F_l ,
          \qquad
          C_{jl} = \sum_n \frac{\varphi_n(\mathbf{x}_j)\,\varphi_n(\mathbf{x}_l)^*}{\lambda_n^0 - \lambda} .
        \end{equation}
  \item Each resonator gives $F_l = Z_l\, w_l$. In matrix form $w = C Z w$, with
        $Z$ diagonal, so a nonzero solution needs
        \begin{equation}
          \boxed{\;\det(I - C Z) = 0\;}
        \end{equation}
\end{enumerate}
With one resonator this is $1 = ZC$. $C_{jl}$ is the deflection at resonator $j$
per unit force at resonator $l$. Each added resonator adds one band and can
couple to patterns that the others cannot.

\paragraph{Where the plane-wave model stops being valid.}
It assumes a uniform plate and point attachments. Cutouts, stiffeners or real
attachment patches need FE. Inside its validity it is nearly free to evaluate, so
it is both the check on FE and a fast design tool.

% ---------------------------------------------------------------------
\section{How detuning moves the edges}

\paragraph{Raising $f_r$.}
$Y = (1 - \lambda/\lambda_r)/\mu$ keeps its value $1/\mu$ at $\lambda = 0$, while
its zero moves to the right, so $Y$ increases at every $\lambda > 0$. At the old
root $Y - P$ is now positive, and since $Y - P$ falls with $\lambda$, the root
moves right: \emph{every coupled root rises with $f_r$}.

\paragraph{Raising $\mu$.}
$Y$ is zero at $\lambda = \lambda_r$ for any $\mu$, so raising $\mu$ turns the
line about that point. Below $\lambda_r$, $Y$ is positive and shrinks, so roots
below $f_r$ go \emph{down}. Above $\lambda_r$, $Y$ is negative and gets closer to
zero, so roots above $f_r$ go \emph{up}. A heavier resonator pushes roots
\emph{away from} $f_r$, which is level repulsion. Check at $f_r = 1200\,\mathrm{Hz}$,
$\mu = 0.25 \to 0.5$: band 3 at $X$ goes from $1653$ to $1842\,\mathrm{Hz}$, band 1
at $M$ from $1032$ to $927\,\mathrm{Hz}$.

\begin{center}
\begin{tabular}{|l|l|l|l|}
\hline
$f_r$ [Hz] & $M$ band 1 [Hz] & $\Gamma$ band 2 [Hz] & $\Gamma$ band 2 $/\, f_r$ \\
\hline
300 & 297 & 334 & 1.11 \\
800 & 748 & 874 & 1.09 \\
1200 & 1031 & 1273 & 1.06 \\
1500 & 1188 & 1546 & 1.03 \\
2000 & 1360 & 1940 & 0.97 \\
5000 & 1600 & 2957 & 0.59 \\
200000 & 1648 & 3279 & 0.02 \\
\hline
\end{tabular}
\end{center}

\paragraph{Limits.}
\begin{itemize}
  \item $f_r \to 0$: $Y$ falls very steeply, so the root sits at
        $\lambda \approx \lambda_r$, and band 1 at $M$ follows $f_r$.
  \item $f_r \to \infty$: $Y$ becomes the horizontal line $1/\mu$, so
        Eq.~\eqref{eq:Y_equals_P} is the rigid-mass condition. The spring is
        infinitely stiff, and band 1 at $M$ saturates at $1648\,\mathrm{Hz}$. The
        four-wave rigid-mass formula $2378/\sqrt{1+4\mu}$ gives
        $1682\,\mathrm{Hz}$; the other waves account for the rest.
\end{itemize}

\paragraph{Three tuning values mark the changes.}
\begin{center}
\begin{tabular}{|l|p{0.62\textwidth}|}
\hline
$f_r$ & What happens \\
\hline
$\approx 1118\,\mathrm{Hz}$ & $\Gamma$ band 2 reaches $f_X$, and the upper edge becomes pinned \\
$\approx 1513\,\mathrm{Hz}$ & $M$ band 1 reaches $f_X$, and the gap closes \\
$\approx 1765\,\mathrm{Hz}$ & $\Gamma$ band 2 equals $f_r$: $S$ reaches $1$, and the plate at $\Gamma$ turns from mass-like to spring-like \\
\hline
\end{tabular}
\end{center}
\paragraph{Why the gap is widest near $1118\,\mathrm{Hz}$.}
Below it, both edges move up with $f_r$, the upper one faster, so the width
grows. Above it, the upper edge is frozen while the lower one keeps rising, so
the width shrinks.

% ---------------------------------------------------------------------
\section{Two resonators per cell: numbers}
\label{sec:two_resonators}

Identical resonators, $f_r = 1200\,\mathrm{Hz}$, plane-wave model on
$\Gamma$--$X$--$M$--$Y$--$\Gamma$ with $15$ $\mathbf{k}$-points per segment. The
model reproduces the one-resonator values to about $1\,\mathrm{Hz}$, but it is not
FE. The path needs $Y$ because a resonator at $(a/2, 0)$ breaks the square
symmetry. The two-resonator gaps lie between bands 2 and 3, because each resonator
adds a band.
\begin{center}
\begin{tabular}{|p{0.23\textwidth}|p{0.15\textwidth}|p{0.18\textwidth}|l|l|}
\hline
Resonators per cell & At $X$ [Hz] & At $Y$ [Hz] & Gap [Hz] & Width \\
\hline
one at $(0,0)$, $\mu = 0.25$ & 835, \textbf{1194}, 1653 & 835, \textbf{1194}, 1653 & 1031--1194 & 163 \\
one at $(0,0)$, $\mu = 0.5$ & 724, \textbf{1194}, 1841 & 724, \textbf{1194}, 1841 & 926--1194 & 269 \\
$(0,0) + (a/2, 0)$, $\mu = 0.25$ each & 835 $\times 2$, 1653 $\times 2$ & 735, 1137, \textbf{1194}, 1913 & 1137--1194 & 58 \\
$(0,0) + (a/2, a/2)$, $\mu = 0.25$ each & 835 $\times 2$, 1653 $\times 2$ & 835 $\times 2$, 1653 $\times 2$ & 1114--1429 & 315 \\
\hline
\end{tabular}
\end{center}

\paragraph{The diagonal arrangement is a smaller lattice.}
Two identical resonators at $(0,0)$ and $(a/2, a/2)$ form a square lattice rotated
by $45^\circ$ with spacing $a/\sqrt{2}$. The first Bragg frequency moves up, from
$X$ at $1194\,\mathrm{Hz}$ to the old $M$ at $2378\,\mathrm{Hz}$. A second resonator
and a smaller cell are therefore the same physics here. The gap then contains
$f_r$, so it is a local-resonance gap.

\paragraph{Extra question: why is the $(0,0)+(a/2,0)$ gap so narrow?}
Sorted values at $Y$:
\begin{center}
\begin{tabular}{|l|l|p{0.58\textwidth}|}
\hline
Band at $Y$ & $f$ [Hz] & Mode \\
\hline
1 & 735 & $\cos(\pi y/a)$, both resonators moving together, lower branch \\
2 & \textbf{1137} & the two resonators moving in antiphase, new mode \\
3 & \textbf{1194.4} & $\sin(\pi y/a)$, node at both resonators, bare \\
4 & 1913 & $\cos(\pi y/a)$, both resonators moving together, upper branch \\
\hline
\end{tabular}
\end{center}
$1137\,\mathrm{Hz}$ is the \emph{maximum} of band 2, so it is the new lower edge;
the upper edge is the minimum of band 3, the pinned mode. The gap shrinks because
the second resonator adds a band, and at $Y$ that band lands inside the old gap.

Why it sits at $1137\,\mathrm{Hz}$: both resonators have $y = 0$, so the waves at
$Y$, which vary only in $y$, have the same value at both. If the resonators move
in antiphase, $F_1 = -F_2$, their forces on those waves cancel exactly, and the $Y$
waves stay at rest. The opposite forces do excite waves that alternate in $x$
between the two resonators: $(\pm 2\pi/a, \pm\pi/a)$, at about
$5868\,\mathrm{Hz}$, and higher waves with odd multiples of $2\pi/a$ in $x$. For
such a wave the two resonators contribute with opposite sign and the antiphase
forces add up, giving a factor $2$:
\begin{equation}
  1 = 2\,\mu_\mathrm{eff}\, P_\mathrm{odd},
  \qquad
  P_\mathrm{odd} = \sum_{\text{odd waves}} \frac{\lambda}{\lambda_n^0 - \lambda} .
\end{equation}
Insert $\mu_\mathrm{eff} = \mu/(1-\beta)$ and multiply by $(1-\beta)$:
\begin{equation}
  1 - \beta = 2\mu\, P_\mathrm{odd} .
\end{equation}
All of these waves lie far above $1137\,\mathrm{Hz}$, so $P_\mathrm{odd} > 0$ and
the plate acts like a soft spring, which pulls the mode below $f_r$. The four waves
at $5868\,\mathrm{Hz}$ give $P_\mathrm{odd} \approx 0.16$, and the higher waves
bring it to about $0.2$, so
\begin{equation}
  \beta \approx 1 - 2\cdot 0.25\cdot 0.205 = 0.897 ,
  \qquad
  f \approx 1200\cdot\sqrt{0.897} \approx 1137\,\mathrm{Hz}.
\end{equation}

% ---------------------------------------------------------------------
\section{Check questions with answers}

\paragraph{Question 1.} \emph{As $f_r$ goes up, band 1 at $M$ rises but saturates
near $1650\,\mathrm{Hz}$, well below its bare value of $2378\,\mathrm{Hz}$. Why does
it stop there, and what has the resonator physically become?}

\emph{Answer.} With $f_r \to \infty$ the spring cannot deform, so the resonator is
a rigid mass glued to the plate. Band 1 at $M$ settles at the mass-loaded value:
$1682\,\mathrm{Hz}$ from the four-wave formula, $1648\,\mathrm{Hz}$ with all waves.
A rigid mass can only lower a frequency, so stiffening the spring further changes
nothing.

\paragraph{Question 2.} \emph{A heavier resonator pushes roots away from $f_r$.
Where in the $Y$--$P$ picture do you see this, and what does it mean for a gap that
contains $f_r$ compared with one lying entirely below $f_r$?}

\emph{Answer.} Raising $\mu$ turns $Y$ about its zero at $\lambda = \lambda_r$.
Roots left of that point move left, roots right of it move right.
\begin{center}
\begin{tabular}{|p{0.25\textwidth}|l|l|p{0.17\textwidth}|}
\hline
Tuning & Gap at $\mu = 0.25$ & Gap at $\mu = 0.5$ & What moved \\
\hline
$f_r = 800\,\mathrm{Hz}$ (contains $f_r$) & 748--874\,Hz & 707--936\,Hz & both edges, in opposite directions \\
$f_r = 1200\,\mathrm{Hz}$ (below $f_r$) & 1032--1194\,Hz & 927--1194\,Hz & lower edge only \\
\hline
\end{tabular}
\end{center}
Widening does not make the $1200\,\mathrm{Hz}$ gap a local-resonance gap. A gap that
contains $f_r$ has one edge on each side of the pivot and is a true local-resonance
gap. The $1200\,\mathrm{Hz}$ gap lies entirely below it, and its upper edge is the
sin pattern at $X$, which is not a root of $Y = P$ at all because it has a node at
the resonator. It stays a Bragg gap with a resonance-boosted mass.

\paragraph{Question 3.} \emph{Add a second resonator in each cell at $(a/2, 0)$, same
$f_r$. What happens to the sin pattern at $X$, and would you expect the upper edge to
stay pinned at $1194\,\mathrm{Hz}$?}

\emph{Answer.} At $X$, $\cos(\pi x/a)$ has its antinode at the first resonator and a
node at the second, and $\sin(\pi x/a)$ the reverse. That makes two identical
plate-plus-resonator pairs. The sin pattern does not simply rise; it \emph{splits}
into a lower and an upper branch, like cos, and the $1194.4\,\mathrm{Hz}$ mode is
gone from $X$.

But $X$ is not the only point that matters. At $Y$ the pattern $\sin(\pi y/a)$ is
zero at $y = 0$, where both resonators sit, so it stays at $1194.4\,\mathrm{Hz}$ and
the upper edge stays pinned. Moving the second resonator to the cell centre removes
the pinning at both $X$ and $Y$ (Section~\ref{sec:two_resonators}).
